Modern software development often requires the collection of \emph{telemetry}, typically information gathered from users that reveals how a product is used. This information is invaluable for identifying software defects (e.g.\ via crash or bug reports) and for guiding product development efforts, such as by tracking how and when users interact with deployed features.

However, unrestricted telemetry access can lead to privacy concerns.
For example, a system that records searched URLs may inadvertently expose sensitive user information, especially in cases where a URL is accessed by only a small number of users, such as a webpage related to a specific medical condition.
The need to respect privacy concerns has spurred substantial work in the area of privacy-preserving statistics calculations ~\cite{CCS:BBGLR20, SP:LBVKH23, CCS:BIKMMP17, AC:LLPT23, SP:MWAPR23, ASIACCS:BSHV23, USENIX:BGLLMR23, EPRINT:DPRS23, prio, SCN:AGJOP22, EPRINT:BanCheVen24, PoPETS:MouSarTso24, PoPETS:MPDST25, SP:BBCGI21, EPRINT:BBCGKR24, zhu2020federated, SP:RSWP23, bittau2017prochlo, bassily2017practical, STOC:BasSmi15, bun2019heavy, CCS:QYYKXR16} including real-word deployments \autocite{apple2021exposure} and recent standardization efforts by the IETF \autocite{ietf-ppm-dap-13}.
One particular area that has received significant attention in the private telemetry space is the \emph{private $\numheavyhit$-heavy-hitters problem}~\cite{CCS:QYYKXR16, CCS:BohKer21, SP:BBCGI21, EPRINT:DPRS23, EPRINT:JKKPGS22, CCS:DSQGLH22, USENIX:LiNavTes24, SP:BBCGKR25}, where a data aggregator only learns responses that are reported by at least $\numheavyhit$ clients.
In the previous example, this would allow the service provider to obtain a list of  popular URLs while keeping rarer and thus potentially more sensitive searches private.

In previous work this problem has seen two approaches.
One direction focuses on concretely efficient constructions from highly optimized two-party computation (2PC)   protocols \autocite{\citepoplar,EPRINT:JKKPGS22,SP:BBCGKR25}.
While efficient, these works require multiple  servers, and the strong assumptions that $(1)$ the enlisted servers will not collude, despite being in constant interactive communication, and $(2)$ that it will be possible for a data aggregator to identify an honest third-party that is willing to expend computational resources equal to the data aggregator for little to no benefit.
Some earlier works \autocite{SP:BBCGI21, SP:BBCGKR25} are also impacted by inherent leakage from the MPC functionality, which they reduce using differential privacy \autocite{TCC:DMNS06}.


To address these limitations, In CCS '22 Davidson et al. proposed STAR \cite{\citestar}.
STAR introduced a new paradigm with two non-colluding servers between which no interaction needs to occur. Critically, one server can even be relatively lightweight.
In this model, clients report their inputs in two phases: first, they run a protocol with a \emph{Randomness Server} to compute a \emph{report}.
The Randomness Server has minimal resource requirements: its interaction with each client consumes only small amounts of bandwidth, is usually computationally inexpensive, and results in no increased storage for the server.
After interacting with the Randomness Server, clients transmit their report to an independent \emph{Aggregation Server}, which uses the collected reports to \emph{locally} compute the $\numheavyhit$-heavy-hitters.
\cref{fig:star-model} provides an overview of the STAR model.


\begin{figure}[ht]
	\centering
	\resizebox{0.6\linewidth}{!}{\begin{tikzpicture}[
		chip/.style={
				draw,
				fill=gray!10,
				thick,
				minimum width=1cm,
				minimum height=1cm,
			},
		pin/.style={
				line width=0.7pt,
			},
		server/.style={
				draw,
				fill=black!85,
				rounded corners=2pt,
				minimum width=2.5cm,
				minimum height=1.25cm,
				drop shadow,
			},
		bay/.style={
				fill=gray!20,
				draw=gray!50,
				rounded corners=0.5pt,
				minimum height=0.15cm,
				minimum width=0.3cm,
			},
		light/.style={
				circle,
				minimum size=2pt,
				inner sep=0pt,
			},
		every node/.style={font=\sffamily}
	]

	% --- Client ---
	\node[alice, minimum size=1cm] (user) at (0,4) {Client};

	% --- Randomness Server ---
	\node[chip] (chip) at (-2,0) {};
	\foreach \x in {-0.25, 0, 0.25} {
			\draw[pin] ($(chip.north)+(\x,0)$) -- ++(0,0.3);
			\draw[pin] ($(chip.south)+(\x,0)$) -- ++(0,-0.3);
		}
	\foreach \y in {-0.25, 0, 0.25} {
			\draw[pin] ($(chip.west)+(0,\y)$) -- ++(-0.3,0);
			\draw[pin] ($(chip.east)+(0,\y)$) -- ++(0.3,0);
		}
	\node at (chip.south) [below=1em] {\textbf{Randomness Server}};

	% --- Aggregation Server ---
	\node[server] (srv) at (2,0) {};
	\foreach \x in {-0.9,-0.6,-0.3,0,0.3,0.6} {
			\foreach \y in {0.25, -0.25} {
					\node[bay] at ($(srv.center) + (\x,\y)$) {};
				}
		}
	\foreach \i/\col in {0.25/green, 0/yellow, -0.25/red} {
			\node[light, fill=\col] at ($(srv.east) + (-0.25,\i)$) {};
		}

	\node at (srv.south) [below=0.4em] {\textbf{Aggregation Server}};

	% Arrows from user to chip and server
	\draw[<->, dashed, thick] ($(user.south) + (-0.3,-0.5)$) -- node[midway, left, yshift=15pt, align=center] {(1) Protocol to \\ generate report} ($(chip.north) + (0,0.5)$);
	\draw[->, thick] ($(user.south) + (0.3,-0.5)$) -- node[midway, right, yshift=8pt] {(2) Send report} ($(srv.north) + (0,0.3)$);
	\node[below=2em of srv] [align=center] {(3) Aggregate to \\ compute heavy-hitters};
\end{tikzpicture}
}
	\caption{The STAR model. Note that there is an implicit assymetry in the computational resources of the Randomness Server and the Aggregation Server.}
	\label{fig:star-model}
\end{figure}
%
% Unfortunately, the STAR protocol suffers from substantial privacy leakage: the Aggregation Server learns a pseudonomized version of the \emph{complete frequency histogram} of client inputs.
% Recognizing this issue, a follow-up work called POPSTAR (Li et al., USENIX 2024) reduces the leakage to a partial histogram, where only a subset of values have their frequency leaked.
% Notably, the values that have their frequencies leaked are selected \emph{at random}, with rare values having the same probability of being leaked as those with more reports.
% More critically, both STAR and POPSTAR are extremely vulnerable to an Aggregation Server that colludes with (or poses as) a client.
% An Aggregation Server that makes a \emph{single query} to the Randomness Server
% can identify \emph{all} reports of a target input value $\val$, even if $\val$ is not a heavy-hitter.
% While this ``spying'' attack was acknowledged in prior work, it was explicitly defined to be outside the threat model~\autocite{\citestar, \citepopstar}, despite being a realistic threat to any deployment.
% Continuing the URL example, an Aggregation Server could use this attack to identify all users visiting a site related to pancreatic cancer.
% Such a value is unlikely to be a heavy-hitter, and reveals considerable information about those reporting it.
% Finally, we emphasize that this attack is possible even in the \emph{honest-but-curious} setting, as it only requires the Aggregation Server to learn a client's view during the report generation phase.
%
%
% \subsection{Our Contributions: Guaranteed Protection for Rare Values} The weaknesses described above motivate the need for a $t$-heavy-hitters protocol with new security guarantees.
% Ideally, a protocol would both \emph{guarantee} protection for rare values, granting total privacy to those reported infrequently, and have hardened security against spying attacks.
%
% We note here that any system in which the Aggregation Server can freely query the Randomness Server and aggregate shares must allow some sort of spying.
% For a value of interest $v$ the Aggregation Server can create $t - 1$ shares for $v$ via the Randomness Server, and then can combine those shares with every submitted report to tell if it contains $v$.
% While such an attack is inevitable, we can still hope to force the Aggregation Server to make a large number of queries to the Randomness Server in order to launch the attack, rather than just a single query per-value as in STAR and POPSTAR.
% Such a limitation would allow deployed systems to more effectively detect a misbehaving Aggregation Server.
%
% Given this, our contributions are as follows:
% \begin{enumerate}
% 	\item We construct the first $t$-heavy-hitters protocol in the STAR model that guarantees perfect security for rare values, i.e.\ those reported under $k$ times for a well defined privacy threshold $k < t$.
%
% 	\item We extend our protocol to construct the first STAR-model $t$-heavy-hitters protocol with \emph{spying resistance}, requiring the optimal $k - 1$ queries per-value to launch a spying attack.
% \end{enumerate}
% We prove our protocol secure in the honest-but-curious setting, and show that it is efficient enough to be deployed today.
% Report generation is under three seconds, and aggegation runs in approximately 1 hour for $100,000$ reports.
% A comparison between this and prior work can be seen in \cref{fig:comparison}.
%
% The constructions of STAR and POPSTAR are fundamentally based on secret sharing, and the main novelties of this paper are in the construction, application, and optimization of secret sharing systems.
% We achieve perfect security for rare values by using an optimized implementation of \emph{multi-dealer secret sharing} \cite{\citeairtag}, and spying resistance via the construction of a novel \emph{oblivious share sampling} protocol, a 2-party protocol that allows one party to sample a secret share of a value without learning the full sharing.
%
% \begin{figure}
% 	\centering
% 	\begin{tabular}{ |c|c|c|c| }
% 		\hline
% 		\makecell{Scheme}              &
% 		\makecell{Leakage}             &
% 		\makecell{Spying\\ Resistance} &
% 		\makecell{Agg\\ Runtime}                                                             \\
% 		\hline
% 		STAR                           & \makecell{full\\ histogram}    & \ding{55} & 2 sec  \\
% 		\hline
% 		POPSTAR                        & \makecell{partial\\ histogram} & \ding{55} & 10 sec \\
% 		\hline
% 		\textbf{This Work}             & \makecell{privacy\\ threshold} & \ding{51} & 1 hour \\
% 		\hline
% 	\end{tabular}
% 	\caption{A comparison between this and prior work. Aggregation runtimes are over $100,000$ reports.}
% 	\label{fig:comparison}
% \end{figure}
%
%
%
%
%
%
